\documentclass[11pt]{article}
\usepackage[margin=1in]{geometry}
\usepackage{amsmath,amssymb}
\usepackage{array}
\newcolumntype{L}[1]{>{\raggedright\arraybackslash}p{#1}}
\usepackage{xcolor}
\usepackage[colorlinks=true,linkcolor=blue!60!black,citecolor=blue!60!black,urlcolor=blue!60!black]{hyperref}

\newcommand{\Klate}{K_{\mathrm{late}}}
\newcommand{\Svar}{\Sigma^2}

\title{The complexity arrow tracks spectral rigidity --- as a trend, not a law:\\
a pre-registered falsifier test on Krylov chains,\\ random-matrix ensembles, and the DSSYK chord basis}
\author{Ariadne Vale, Tobin Masala, and agy-science\\[2pt]
\small Executable Positive Geometry\\[-2pt]
\small \texttt{epg@subnet345.com}}
\date{29 July 2026}

\begin{document}
\maketitle

\begin{abstract}
Krylov (spread) complexity grows on a semi-infinite chain whose position
operator is positive with a boundary at zero --- a structure that in
double-scaled SYK (DSSYK) is not an analogy but an identity: the thermofield
double's Krylov basis \emph{is} the chord-number basis, so complexity equals
the expected chord number~\cite{rsss}. We ask what property of a Hamiltonian's
\emph{spectrum} controls this complexity arrow, and we ask it as a
pre-registered falsifier test. Across seven ensemble families --- GOE, Poisson,
Anderson chains, $\beta$-Hermite tridiagonal, Rosenzweig--Porter, banded, and
the deterministic DSSYK chord chain --- with spectral rigidity measured from
independently unfolded finite spectra (never from the Jacobi data driving the
complexity), we find: (1)~no counterexample to rigidity--arrow co-variation,
including the ensemble designed to break it; (2)~a forced refinement --- the
DSSYK chord chain is \emph{picket-fence} rigid (gap ratio $r=0.968$,
integrable-type, not random-matrix) yet carries the strongest arrow, so the
arrow tracks \emph{rigidity}, not chaos; chaos is one route to rigidity;
(3)~quantitatively, the trend is real (Spearman $\rho=-0.765$,
$p=3.4\times10^{-5}$, pre-registered) but \emph{universality fails exactly as
its pre-registered falsifier anticipated}: at matched rigidity, ensemble
families branch by an order of magnitude in complexity, and the branch
variable is eigenvector localization (inverse participation ratio). The honest
conclusion: \emph{spectral rigidity predicts the complexity arrow as a trend
but is not sufficient; eigenvector delocalization is the dominant additional
variable.}
Every claim was pre-registered with its falsifier before running; one of our
own prior headline claims (``the arrow ignites from the chaotic spectrum'')
is narrowed by these results, and we say so.
\end{abstract}

\section{Setup: the arrow lives on a one-sided chain}

The Lanczos recursion maps any Hamiltonian and seed state to a semi-infinite
tridiagonal chain; spread complexity $K(t)=\langle \hat n\rangle$ is the
wavepacket's mean position, a manifestly positive operator with a boundary at
$n=0$. In DSSYK this chain is physical: the chord
basis~\cite{lin,berkooz} makes the Hamiltonian tridiagonal with hopping
$b_n=\sqrt{(1-q^n)/(1-q)}$, and the thermofield double's Krylov basis
coincides with the chord-number basis exactly, at any $q$~\cite{rsss} ---
complexity \emph{is} chord number, and (in the triple-scaling limit) bulk
wormhole length; at finite $N$ the chord basis is an extrapolation and
complexity saturates~\cite{bmnw}. Our exact chain lab confirms the registered
expectations: $K(t)=b_1^2t^2$ at small $t$ (to $<2\times10^{-5}$); ballistic
pre-boundary transport at every $q$ with velocity set by the hopping plateau
$b_\infty=1/\sqrt{1-q}$; a super-ballistic ($K\!\sim\!t^2$) contrast in the
$q\to1$ oscillator limit; all readouts strictly inside the pre-reflection
window, so no finite-cutoff artifact is claimed as physics.

\section{The falsifier design}

The Lanczos coefficients determine both the spectral density and $K(t)$, so
``complexity relates to the spectrum'' is trivially true. To make the question
falsifiable we required: (i)~spectral rigidity measured from
\emph{independently unfolded finite spectra} --- gap ratios $r$ and number
variance $\Svar(L)$ --- never reconstructed from the same Jacobi data that
evolves $K(t)$; (ii)~ensembles chosen to break the hypothesis if it can be
broken, in particular $\beta$-Hermite tridiagonal matrices (disordered Lanczos
coefficients with exactly random-matrix spectra) and Rosenzweig--Porter
interpolation through its transition; (iii)~verdicts pre-registered before
each run, with the falsifier stated. An independent analytical reviewer
certified the design non-circular before any code ran.

\section{Stage one: no counterexample, one refinement}

Across GOE, Poisson, Anderson, $\beta$-Hermite, Rosenzweig--Porter, banded,
and DSSYK-chord ensembles (two independent executions; dimensions 64--128, up
to 64 disorder samples), no clean bridge-killer appeared: rigid spectra came
with delocalized chains and strong late complexity, Poisson-like spectra with
localized chains and weak complexity, intermediate ensembles intermediate on
both axes. The designed killer survived: $\beta$-Hermite's coefficient
disorder is relatively weak and both its columns stay random-matrix-like.

The refinement came from our own arena. The DSSYK chord chain's spectrum has
$r=0.968$ with low number variance: near-equally-spaced levels ---
\emph{picket-fence} rigidity, characteristic of integrable spectra like the
oscillator, categorically distinct from random-matrix level repulsion
($r\approx0.53$--$0.60$) --- yet it carries the strongest arrow in the
comparison. The classifier was patched to separate picket-fence from
random-matrix rigidity, and the bridge statement was narrowed accordingly:
\begin{quote}
\emph{The complexity arrow co-varies with spectral rigidity (low number
variance), not with chaos per se. Chaos is one route to rigidity; the
deterministic chord chain is another.}
\end{quote}
This narrows our own earlier claim (``the arrow ignites from the chaotic
spectrum''): chaotic should have been \emph{rigid}. The chord chain, being a
single deterministic instance, is kept in all comparisons as a labeled $n=1$
structured control, not as an ensemble member.

\section{Stage two: the trend is real; universality fails, informatively}

We then pre-registered a quantitative test over 25 configurations --- 22
ensemble scatter points plus 3 deterministic chord controls --- spanning the
seven families (parameter sweeps in disorder strength, $\beta$, coupling,
bandwidth, and $q$), with all measurement code identical to stage one:

\textbf{Q-1 (trend) --- confirmed.} Normalized late complexity
$\Klate/\mathrm{chain}$ versus $\Svar(L{=}4)$ is monotone-decreasing across
the ensemble scatter: Spearman $\rho=-0.765$, $p=3.4\times10^{-5}$ at
dimension 128 with 32 samples per point (independent smoke run at lower power:
$\rho=-0.743$).

\textbf{Q-2 (universality) --- refuted, exactly as pre-registered.} If
rigidity \emph{sufficed}, points from different families at matched $\Svar$
would agree. They do not, in any populated rigidity bin: e.g.\ at
$\Svar\approx3.5$--$4.0$, an Anderson chain ($W{=}5$) gives
$\Klate/\mathrm{chain}=0.032$ while Rosenzweig--Porter ($c{=}0.5$) gives
$0.481$ --- a factor $\sim$15 at matched rigidity, far beyond combined
uncertainties. The pre-registered fallback identified the branch variable in
advance: eigenvector localization. The recorded inverse participation ratios
at that matched point are $0.328$ (localized) versus $0.067$ (delocalized).

\textbf{Q-3 (control) --- consistent.} The deterministic chord points
($q=0.05/0.5/0.9$) sit in the rigid--strong region of the scatter: at the
registered paper-grade dimension, $\Svar = 2.86/2.28/1.44$ with
$\Klate/\mathrm{chain} = 0.546/0.570/0.408$ and gap ratios up to $0.986$ ---
strong complexity at modest number variance, consistent with the ensemble
trend. (Their number variance is dimension-dependent through the unfolding of
a deterministic spectrum; we therefore make no superlative rigidity claim for
them, and rest their picket-fence character on the gap ratio.)

\section{Conclusion, stated honestly}

\begin{quote}
\emph{Spectral rigidity predicts the complexity arrow as a trend
($\rho\approx-0.77$) but is not sufficient: there is no universal
$\Klate(\Svar)$ curve. Eigenvector delocalization is the dominant additional
variable; other factors may remain.}
\end{quote}

Three scope statements. First, these are finite-dimensional numerical results
(shadows of the type-III$_1$/large-$N$ statements they gesture at); a
falsifier hunt that finds no counterexample is support, not proof. Second, the
DSSYK chord chain enters as a deterministic labeled control; its picket-fence
rigidity is an observation about one physically central instance, not an
ensemble statement. Third, the underlying conjecture that motivated this study
--- that the one-sided chord chain's positivity ($\hat n\ge0$ with a boundary)
is the finite shadow of a half-sided modular/arrow-of-time structure in DSSYK
--- remains exactly that, a conjecture; our literature verification found the
positivity-as-arrow fusion unoccupied (positivity is used thermodynamically
in~\cite{bmp}, modular structure algebraically in~\cite{xu}), and nothing here
proves it.

\section*{Methods --- pre-registration and three-agent verification}

Every stage followed the same discipline: predictions and falsifiers
registered in the public repository before running; measurement code frozen
across stages (the quantitative stage calls the stage-one measurement function
verbatim, sweeping only parameters); spectra unfolded independently of the
chain data; two independent executions (author and clean-room, the latter from
the formulas alone) for every load-bearing number; and an analytical
reviewer's audit at each gate (non-circularity of the design; the
$b_n$ symmetrization; the picket-fence reading; the final conclusion's
wording). Two claims of ours died in the process --- the universality of
$\Klate(\Svar)$, and the ``chaotic spectrum'' phrasing of our own earlier
capstone --- and both deaths are reported as results. Produced by three AI
research agents on distinct runtimes (Anthropic, OpenAI, and a third
analytical reviewer) under human direction; every number carries a commit
hash and a persisted run manifest.

\begin{thebibliography}{9}
\bibitem{rsss} E.~Rabinovici, A.~S\'anchez-Garrido, R.~Shir, J.~Sonner,
  \emph{A bulk manifestation of Krylov complexity}, JHEP \textbf{08} (2023)
  213, \href{https://arxiv.org/abs/2305.04355}{arXiv:2305.04355}; see also
  \href{https://arxiv.org/abs/2412.15318}{arXiv:2412.15318}.
\bibitem{lin} H.~W.~Lin, \emph{The bulk Hilbert space of double scaled SYK},
  \href{https://arxiv.org/abs/2208.07032}{arXiv:2208.07032}.
\bibitem{berkooz} M.~Berkooz, M.~Isachenkov, V.~Narovlansky, G.~Torrents,
  \emph{Towards a full solution of the large N double-scaled SYK model},
  \href{https://arxiv.org/abs/1811.02584}{arXiv:1811.02584}.
\bibitem{bmnw} V.~Balasubramanian, J.~M.~Magan, P.~Nandi, Q.~Wu,
  \emph{Spread complexity and the saturation of wormhole size},
  \href{https://arxiv.org/abs/2412.02038}{arXiv:2412.02038}.
\bibitem{parker} D.~E.~Parker, X.~Cao, A.~Avdoshkin, T.~Scaffidi, E.~Altman,
  \emph{A Universal Operator Growth Hypothesis},
  \href{https://arxiv.org/abs/1812.08657}{arXiv:1812.08657}.
\bibitem{bmp} A.~Blommaert, T.~G.~Mertens, J.~Papalini, \emph{The dilaton
  gravity hologram of double-scaled SYK},
  \href{https://arxiv.org/abs/2404.03535}{arXiv:2404.03535}.
\bibitem{xu} J.~Xu, \emph{Von Neumann Algebras in Double-Scaled SYK},
  \href{https://arxiv.org/abs/2403.09021}{arXiv:2403.09021}.
\bibitem{heller} M.~P.~Heller, F.~Ori, J.~Papalini, T.~Schuhmann,
  M.-T.~Wang, \emph{De Sitter holographic complexity from Krylov complexity
  in DSSYK}, \href{https://arxiv.org/abs/2510.13986}{arXiv:2510.13986}.
\end{thebibliography}

\vspace{1em}
\noindent\rule{\linewidth}{0.4pt}\\[2pt]
{\footnotesize \emph{Acknowledgements.} Produced by a swarm of AI research
agents --- OpenAI, Anthropic, and a third analytical-reviewer agent --- under
human direction, each independently checking the others before a claim was
allowed to stand. Code, run manifests, and the pre-registration records are
public; refuted predictions are reported beside confirmations.}

\end{document}
